
\bta{2006乙A}


\begin{enumerate}
	%\renewcommand{\labelenumi}{\arabic{enumi}.}
	% A(\Alph) a(\alph) I(\Roman) i(\roman) 1(\arabic)
	%设定全局标号series=example	%引用全局变量resume=example
	%[topsep=-0.3em,parsep=-0.3em,itemsep=-0.3em,partopsep=-0.3em]
	%可使用leftmargin调整列表环境左边的空白长度 [leftmargin=0em]
	\item
（$30'$）粒子以能量$E$入射方势垒，$V(x)=\begin{cases}V_0>0,&0\le x\le a\\0,&x<0,x>0\end{cases}$。设能量$E>V_0$，求透射系数$T$。

\banswer{
	令
	\[
		k=\frac{\sqrt{2mE}}{\hbar},\qquad
		q=\frac{\sqrt{2m(E-V_0)}}{\hbar}.
	\]
	三区波函数可分别写为 $e^{ikx}+r e^{-ikx}$、$Ae^{iqx}+Be^{-iqx}$、$t e^{ikx}$。在 $x=0,a$ 连续匹配波函数及其一阶导数，消去 $A,B,r$ 后得到
	\[
		\boxed{T=|t|^2=
		\left[1+\frac{V_0^2}{4E(E-V_0)}\sin^2(qa)\right]^{-1}}.
	\]
	当 $qa=n\pi$ 时出现共振透射，$T=1$。
}


\item 
（$30'$）自旋为$\frac{1}{2}$的粒子置于势场$V(x)$中，$V(x)=\begin{cases}0,&0\le x<a\\ \infty,&x<0,x>a\end{cases}$。设粒子所处状态为$\psi(x,s_z)=\sqrt{\frac{5}{18}}\varphi_3(x)\begin{pmatrix}1\\-i\end{pmatrix}+\sqrt{\frac{2}{9}}\varphi_5(x)\begin{pmatrix}1\\ i\end{pmatrix}$，其中$\varphi_n(x)$为系统空间部分的第$n$个能量本征函数（已归一化）。求能量的可测值及相应的取值概率。


\banswer{
	空间能级为
	\[
		E_n=\frac{n^2\pi^2\hbar^2}{2ma^2}.
	\]
	两项的空间波函数正交，且两个自旋量的模平方均为
	$|(1,\mp i)^T|^2=2$。因此
	\[
		P(E_3)=\frac{5}{18}\times2=\frac59,
		\qquad
		P(E_5)=\frac29\times2=\frac49.
	\]
	故能量可测值及概率为
	\[
		\boxed{E_3=\frac{9\pi^2\hbar^2}{2ma^2},\ P_3=\frac59;\qquad
		E_5=\frac{25\pi^2\hbar^2}{2ma^2},\ P_5=\frac49}.
	\]
}


\item 
用不确定度关系估算一维谐振子的基态能量。
\banswer{
	设基态中 $\langle x\rangle=\langle p\rangle=0$，记 $d=\Delta x$。由 $\Delta p\geq\hbar/(2d)$，
	\[
		E\simeq\frac{(\Delta p)^2}{2m}+\frac12m\omega^2(\Delta x)^2
		\geq\frac{\hbar^2}{8md^2}+\frac12m\omega^2d^2.
	\]
	右端对 $d$ 取极小，得到
	\[
		d^2=\frac{\hbar}{2m\omega},
		\qquad
		\boxed{E_0\simeq\frac12\hbar\omega}.
	\]
	本题中 Gaussian 基态恰好饱和不确定关系，因此该估计等于精确基态能量。
}


\newpage
\item 
（$30'$）各向同性的三维谐振子哈密顿算符为$\hat{H}^0=-\frac{\hbar^2}{2m}\nabla^2+\frac{m\omega^2}{2}r^2$。加上微扰$\hat{H}'=\lambda(xy+yz+zx)$后，求第一激发态的一级能量修正。

\banswer{
	第一激发态的简并基可取
	$|100\rangle,|010\rangle,|001\rangle$。由
	$\langle1|x|0\rangle=\sqrt{\hbar/(2m\omega)}$，在此子空间中
	\[
		W=\frac{\lambda\hbar}{2m\omega}
		\begin{pmatrix}0&1&1\\1&0&1\\1&1&0\end{pmatrix}.
	\]
	该矩阵的本征值为 $2,-1,-1$。因此一级能量修正为
	\[
		\boxed{E_s^{(1)}=\frac{\lambda\hbar}{m\omega}},
		\qquad
		\boxed{E_d^{(1)}=-\frac{\lambda\hbar}{2m\omega}\quad(\text{二重简并})}.
	\]
	非简并组合为 $(|100\rangle+|010\rangle+|001\rangle)/\sqrt3$；与它正交的任意两个归一化组合张成二重简并子空间。无扰第一激发能量为 $5\hbar\omega/2$。
}


\item 
（$30'$）自旋为$\frac{1}{2}$，磁矩为$\mu$，电荷为零的粒子置于磁场中。$t=0$是磁场为$\vec{B}_0=(0,0,B_0)$，粒子处于$\hat{\sigma}_z$的本征值为$-1$的本征态$\begin{pmatrix}0\\1\end{pmatrix}$。设在$t>0$时，再加上弱磁场$\vec{B}_1=(B_1,0,0)$，求$t>0$时的波函数，以及测到自旋反转的概率。

\banswer{
	取磁矩算符为 $\vec\mu=\mu\vec\sigma$，则
	\[
		H=-\mu(B_0\sigma_z+B_1\sigma_x),
		\qquad B=\sqrt{B_0^2+B_1^2}.
	\]
	时间演化算符为
	\[
		U(t)=\cos\frac{\mu Bt}{\hbar},I
		+i\frac{B_0\sigma_z+B_1\sigma_x}{B}
		\sin\frac{\mu Bt}{\hbar}.
	\]
	作用于初态 $(0,1)^T$ 得
	\[
		\boxed{\psi(t)=
		\begin{pmatrix}
		i\dfrac{B_1}{B}\sin(\mu Bt/\hbar)\\[5pt]
		\cos(\mu Bt/\hbar)-i\dfrac{B_0}{B}\sin(\mu Bt/\hbar)
		\end{pmatrix}}.
	\]
	测得自旋反转为 $\sigma_z=+1$ 的概率为
	\[
		\boxed{P_{-\to+}(t)=\frac{B_1^2}{B_0^2+B_1^2}
		\sin^2\!\left(\frac{\mu t}{\hbar}\sqrt{B_0^2+B_1^2}\right)}.
	\]
	在 $B_1\ll B_0$ 时，$P_{-\to+}\simeq(B_1/B_0)^2\sin^2(\mu B_0t/\hbar)$。
}


	
\end{enumerate}
